\subsection{مبانی نظری و فرمول‌بندی ریاضی}
فرآیند فشرده‌سازی در این ماژول طی مراحل ریاضی زیر انجام می‌پذیرد:
\begin{enumerate}
    \item \textbf{گردآوری فعال‌سازی‌ها:} با عبور دسته‌های کالیبراسیون از مدل، ماتریس نمونه فعال‌سازی ورودی لایه $\mathbf{X} \in \mathbb{R}^{N \times d_{\text{in}}}$ با میانگین $\boldsymbol{\mu} \in \mathbb{R}^{d_{\text{in}}}$ جمع‌آوری می‌شود.
    \item \textbf{محاسبه کوواریانس تجربی:} ماتریس کوواریانس نمونه‌ها به صورت زیر محاسبه می‌گردد:
    \begin{equation}
    \boldsymbol{\Sigma} = \frac{1}{N - 1} (\mathbf{X} - \boldsymbol{\mu})^\top (\mathbf{X} - \boldsymbol{\mu}) \in \mathbb{R}^{d_{\text{in}} \times d_{\text{in}}}
    \end{equation}
    \item \textbf{تجزیه ویژه (\lr{Eigendecomposition}):} ماتریس کوواریانس به پایه‌های ارتوگونال تجزیه می‌شود:
    \begin{equation}
    \boldsymbol{\Sigma} = \mathbf{V} \boldsymbol{\Lambda} \mathbf{V}^\top
    \end{equation}
    که در آن مقادیر ویژه $\lambda_1 \ge \lambda_2 \ge \dots \ge \lambda_{d_{\text{in}}} \ge 0$ نشان‌دهنده میزان پراکندگی در امتداد بردارهای ویژه متناظر در ستون‌های $\mathbf{V}$ هستند.
    \item \textbf{تعیین بعد برش بر مبنای واریانس تجمعی:} با تعیین نسبت تنک‌سازی (\lr{sparsity})، بعد موثر $k$ بر اساس سهم واریانس تجمعی حفظ‌شده انتخاب می‌شود:
    \begin{equation}
    k = \min \left\{ r \;\middle|\; \frac{\sum_{i=1}^r \lambda_i}{\sum_{i=1}^{d_{\text{in}}} \lambda_i} \ge 1 - \text{sparsity} \right\}
    \end{equation}
    \item \textbf{تبدیل و نگاشت وزن‌ها:} ماتریس وزن پایه $\mathbf{W} \in \mathbb{R}^{d_{\text{out}} \times d_{\text{in}}}$ با تصویر روی $k$ مؤلفه نخست به صورت متعامد دگرگون و بازسازی می‌شود:
    \begin{equation}
    \mathbf{W}_{\text{sliced}} = (\mathbf{W} \mathbf{V}) \mathbf{P}_k \mathbf{P}_k^\top \mathbf{V}^\top
    \end{equation}
    که در آن $\mathbf{P}_k \in \mathbb{R}^{d_{\text{in}} \times k}$ ماتریس تصویرگر روی $k$ بعد برتر است.
    \item \textbf{ادغام لایه با آداپتور کم‌مرتبه (\lr{LoRA}):} خروجی نهایی لایه منجمدشده با تصحیح کم‌مرتبه محاسبه می‌شود:
    \begin{equation}
    \mathbf{y} = (\mathbf{x} \mathbf{W}_{\text{sliced}}^\top + \mathbf{b}) + \frac{\alpha}{r} (\mathbf{x} \mathbf{A}^\top) \mathbf{B}^\top
    \end{equation}
\end{enumerate}
